\documentclass[11pt,a4paper]{article}
\usepackage{fontspec}
\usepackage{kotex}
\usepackage{geometry}
\usepackage{graphicx}
\usepackage{xcolor}
\usepackage{tcolorbox}
\usepackage{booktabs}
\usepackage{array}
\usepackage{longtable}
\usepackage{colortbl}
\usepackage{fancyhdr}
\usepackage{titlesec}
\usepackage{hyperref}
\usepackage{emoji}
\usepackage{amsmath}
\usepackage{amssymb}
\usepackage{enumitem}

\geometry{left=25mm, right=25mm, top=30mm, bottom=30mm}

\setmainfont{Noto Sans CJK KR}
\setsansfont{Noto Sans CJK KR}
\setmonofont{Noto Sans CJK KR}
\setemojifont{Noto Color Emoji}

\definecolor{primary}{HTML}{1976D2}
\definecolor{darkbrown}{HTML}{3E2723}
\definecolor{mediumbrown}{HTML}{5D4037}
\definecolor{lightbrown}{HTML}{8D6E63}
\definecolor{cream}{HTML}{FFF3E0}
\definecolor{lightgray}{HTML}{F5F5F5}
\definecolor{tableheader}{HTML}{E8E8E8}

\titleformat{\section}
  {\Large\bfseries\color{primary}}
  {\thesection.}{0.5em}{}
\titleformat{\subsection}
  {\large\bfseries\color{darkbrown}}
  {\thesubsection}{0.5em}{}
\titleformat{\subsubsection}
  {\normalsize\bfseries\color{mediumbrown}}
  {\thesubsubsection}{0.5em}{}

\renewcommand{\contentsname}{목차}
\renewcommand{\figurename}{그림}
\renewcommand{\tablename}{표}

\tcbset{
  tipbox/.style={colback=cream,colframe=primary,boxrule=1.5pt,arc=3mm,
    left=5mm,right=5mm,top=3mm,bottom=3mm,fonttitle=\bfseries\color{primary}},
  formulabox/.style={colback=lightgray,colframe=lightgray,boxrule=0pt,arc=2mm,
    left=5mm,right=5mm,top=3mm,bottom=3mm},
  summarybox/.style={colback=white,colframe=primary,boxrule=2pt,arc=4mm,
    left=5mm,right=5mm,top=5mm,bottom=5mm}
}

\hypersetup{colorlinks=true,linkcolor=primary,urlcolor=primary,bookmarks=true,unicode=true}


\pagestyle{fancy}
\fancyhf{}
\fancyhead[R]{\small\color{lightbrown}마이크로비트 Grove I2C 연결도 이론해설서 v1.0.0.0}
\fancyfoot[C]{\thepage}
\renewcommand{\headrulewidth}{0.4pt}

\begin{document}

\begin{titlepage}
\centering
\vspace*{3cm}
{\Huge\bfseries\color{primary}\emoji{electric-plug} 마이크로비트 Grove I2C 연결도}\\[0.5cm]
{\Huge\bfseries\color{primary} 이론해설서}
\\[1.5cm]
{\Large 시뮬레이터 버전: v1.0.3.0}
\\[2cm]
\begin{tcolorbox}[summarybox]
\centering
{\large 이 이론해설서는 마이크로비트 Grove I2C 연결도이 다루는 핵심 개념과 원리를 설명합니다.}
\end{tcolorbox}
\vfill
{\large 사물인터넷과 빅데이터}\\[0.3cm]
{\large 청주교육대학교 컴퓨터교육과}
\vspace{1cm}
\end{titlepage}

\section*{1. 이 문서의 목적}

이 이론해설서는 \textbf{마이크로비트 Grove I2C 연결도}이 다루는 핵심 개념과 원리를 설명합니다. 시뮬레이터를 사용하기 전에 배경 지식을 쌓거나, 사용 후 깊이 있는 이해를 위해 참조합니다.


\begin{tcolorbox}[summarybox]
\textbf{핵심 주제}: I2C 버스와 물리적 연결

\end{tcolorbox}


\section*{2. 핵심 개념}

\begin{center}
\renewcommand{\arraystretch}{1.4}
\begin{tabular}{|p{3.7cm}|p{11.3cm}|}
\hline
\rowcolor{tableheader}
\textbf{개념} & \textbf{설명} \\
\hline
\textbf{I2C 버스 토폴로지} & 하나의 버스 라인에 여러 장치를 병렬로 연결하는 구조 \\
\hline
\textbf{풀업 저항} & I2C 라인을 HIGH 상태로 유지하기 위한 저항 (보통 4.7kΩ) \\
\hline
\textbf{주소 충돌} & 동일한 I2C 주소를 가진 장치가 2개 이상이면 통신 불가 \\
\hline
\textbf{SDA/SCL} & SDA(Serial Data): 데이터 전송선, SCL(Serial Clock): 동기화 클럭선 \\
\hline
\textbf{배선 표준} & Grove 4핀: VCC(전원), GND(접지), SDA(데이터), SCL(클럭) \\
\hline
\end{tabular}
\end{center}


\section*{3. 원리 상세}

\subsection*{I2C 버스의 물리적 구조}

I2C 버스는 \textbf{버스 토폴로지}를 사용합니다. 하나의 SDA/SCL 선에 모든 센서가 병렬로 연결됩니다.


\textbf{왜 2개 선만으로 충분한가?}

$\bullet$ SCL(클럭): 마스터가 "지금 읽어!"라고 타이밍을 알려줌

$\bullet$ SDA(데이터): 실제 데이터(0과 1)가 전송되는 선

$\bullet$ 각 센서는 고유 주소가 있어 자기 차례에만 응답


\textbf{주소 체계 (7비트):}

$\bullet$ 이론적 최대: 128개 장치 (0x00\textasciitilde{}0x7F)

$\bullet$ 실제 사용: 16개 예약 \textrightarrow{} 112개 가용

$\bullet$ Grove 시스템: 보통 6\textasciitilde{}8개 센서 연결


\textbf{배선 규칙:}

1. 모든 장치의 SDA는 하나로, SCL도 하나로 연결

2. VCC(3.3V)와 GND는 각 장치에 공급

3. 풀업 저항은 버스 전체에 1세트만 필요



\section*{4. 실생활 응용 사례}

\textbf{사례 1}: 자동차: ECU 간 통신에 I2C/CAN 버스 사용 (수십 개 센서 연결)


\textbf{사례 2}: 로봇: 모터·센서·디스플레이를 I2C로 연결하여 제어


\textbf{사례 3}: 가전제품: TV 리모컨 수신, 전자레인지 디스플레이 등에 I2C 활용


\textbf{사례 4}: 의료 기기: 심박·혈압·산소포화도 센서를 I2C로 통합



\section*{5. 더 알아보기}

$\bullet$ \textbf{SPI 통신}: I2C보다 빠르지만 선이 더 필요한 대안 프로토콜

$\bullet$ \textbf{UART}: 1:1 직렬 통신 방식 (RS-232)

$\bullet$ \textbf{주소 멀티플렉싱}: 같은 주소 센서를 I2C 멀티플렉서로 여러 개 연결하는 기법

$\bullet$ \textbf{프로토콜 분석기}: 실제 I2C 통신 신호를 오실로스코프로 확인하는 방법



\section*{6. 핵심 용어 사전}

\begin{center}
\renewcommand{\arraystretch}{1.4}
\begin{tabular}{|p{3.0cm}|p{3.0cm}|p{9.1cm}|}
\hline
\rowcolor{tableheader}
\textbf{용어} & \textbf{학술 용어} & \textbf{쉬운 설명} \\
\hline
I2C 버스 토폴로지 & I2C 버스 토폴로지 & 하나의 버스 라인에 여러 장치를 병렬로 연결하는 구조 \\
\hline
풀업 저항 & — & I2C 라인을 HIGH 상태로 유지하기 위한 저항 (보통 4 \\
\hline
주소 충돌 & — & 동일한 I2C 주소를 가진 장치가 2개 이상이면 통신 불가 \\
\hline
SDA/SCL & SDA/SCL & SDA(Serial Data): 데이터 전송선, SCL(Serial Cl \\
\hline
배선 표준 & — & Grove 4핀: VCC(전원), GND(접지), SDA(데이터), SC \\
\hline
\end{tabular}
\end{center}


\section*{7. 참고자료}

\begin{center}
\renewcommand{\arraystretch}{1.4}
\begin{tabular}{|p{1.5cm}|p{7.3cm}|p{6.2cm}|}
\hline
\rowcolor{tableheader}
\textbf{\#} & \textbf{자료명} & \textbf{비고} \\
\hline
1 & 마이크로비트GroveI2C연결도\_퀵가이드 & 소프트웨어 사용 중 빠른 참조 \\
\hline
2 & 마이크로비트GroveI2C연결도\_사용설명서 & 전체 기능 상세 \\
\hline
3 & 마이크로비트GroveI2C연결도\_활용가이드 & 수업 적용 \\
\hline
4 & 마이크로비트GroveI2C연결도\_개발참조 & 기술 사양 \\
\hline
\end{tabular}
\end{center}


\vfill
\begin{center}
\small\color{lightbrown}
\textit{마이크로비트 Grove I2C 연결도 이론해설서}
\end{center}
\end{document}